\documentclass[11pt,a4paper]{article}

%================================================
% PACKAGES
%================================================
\usepackage[utf8]{inputenc}
\usepackage[T1]{fontenc}
\usepackage[french]{babel}
\usepackage[a4paper,margin=2.2cm]{geometry}

\usepackage{lmodern}
\usepackage{amsmath,amssymb,amsthm,mathtools}
\usepackage{mathrsfs}
\usepackage{xcolor}
\usepackage{hyperref}
\usepackage{fancyhdr}
\usepackage{enumitem}
\usepackage{array}
\usepackage{multicol}
\usepackage{tikz}
\usetikzlibrary{arrows.meta,calc,angles,quotes,intersections}

%================================================
% HYPERREF
%================================================
\hypersetup{
  colorlinks=true,
  linkcolor=black,
  urlcolor=blue,
  pdftitle={Cours complet de géométrie jusqu'à la terminale},
  pdfauthor={}
}

%================================================
% MISE EN PAGE
%================================================
\setlength{\parindent}{0pt}
\setlength{\parskip}{0.45em}

\pagestyle{fancy}
\fancyhf{}
\fancyhead[L]{Cours complet de géométrie}
\fancyhead[R]{Collège $\to$ Terminale}
\fancyfoot[C]{\thepage}

%================================================
% COMMANDES
%================================================
\newcommand{\R}{\mathbb{R}}
\newcommand{\Vect}{\mathrm{Vect}}
\newcommand{\norme}[1]{\left\|#1\right\|}
\newcommand{\ps}[2]{\left\langle #1,#2\right\rangle}
\newcommand{\uu}{\overrightarrow{u}}
\newcommand{\vv}{\overrightarrow{v}}
\newcommand{\ww}{\overrightarrow{w}}
\newcommand{\AB}{\overrightarrow{AB}}
\newcommand{\AC}{\overrightarrow{AC}}
\newcommand{\BC}{\overrightarrow{BC}}

%================================================
% ENVIRONNEMENTS
%================================================
\theoremstyle{definition}
\newtheorem{definition}{Définition}[section]
\newtheorem{exemple}[definition]{Exemple}
\newtheorem{remarque}[definition]{Remarque}
\newtheorem{methode}[definition]{Méthode}

\theoremstyle{plain}
\newtheorem{propriete}[definition]{Propriété}
\newtheorem{proposition}[definition]{Proposition}
\newtheorem{theoreme}[definition]{Théorème}
\newtheorem{corollaire}[definition]{Corollaire}
\newtheorem{lemme}[definition]{Lemme}

%================================================
% TITRE
%================================================
\title{\textbf{Cours complet de géométrie jusqu'à la terminale}\\
\large Géométrie plane et géométrie de l'espace\\
\large Définitions, propriétés, propositions, démonstrations, exemples}
\author{}
\date{}

\begin{document}

\maketitle
\tableofcontents
\newpage

%================================================
\section*{Introduction}
%================================================

Ce cours rassemble les principales notions de géométrie rencontrées :
\begin{itemize}
    \item au \textbf{collège} : figures usuelles, triangles, cercle, symétries, parallélisme, perpendicularité, Thalès, Pythagore, trigonométrie ;
    \item en \textbf{seconde} : géométrie repérée et vecteurs du plan ;
    \item en \textbf{première spécialité} : produit scalaire et géométrie métrique du plan ;
    \item en \textbf{terminale spécialité} : géométrie dans l'espace.
\end{itemize}

L'objectif est de maîtriser :
\begin{itemize}
    \item les définitions ;
    \item les propriétés essentielles ;
    \item les démonstrations classiques ;
    \item les méthodes de résolution.
\end{itemize}

\newpage

%================================================
\section{Cycle 4 : espace et géométrie}
%================================================

\subsection{Objets géométriques de base}

\begin{definition}
Un \textbf{point} est un objet géométrique sans dimension.
\end{definition}

\begin{definition}
Une \textbf{droite} est un ensemble infini de points alignés.
\end{definition}

\begin{definition}
Si $A$ et $B$ sont deux points distincts :
\begin{itemize}
    \item la droite passant par $A$ et $B$ se note $(AB)$ ;
    \item le segment d'extrémités $A$ et $B$ se note $[AB]$ ;
    \item la demi-droite d'origine $A$ passant par $B$ se note $[AB)$.
\end{itemize}
\end{definition}

\begin{definition}
Trois points sont \textbf{alignés} s'ils appartiennent à une même droite.
\end{definition}

\begin{propriete}
Par deux points distincts passe une unique droite.
\end{propriete}

\begin{definition}
La \textbf{distance} entre deux points $A$ et $B$ est la longueur du segment $[AB]$, notée $AB$.
\end{definition}

\begin{definition}
Le \textbf{milieu} du segment $[AB]$ est le point $M$ du segment tel que
\[
AM=MB.
\]
\end{definition}

\begin{proposition}
Un segment possède un unique milieu.
\end{proposition}

\begin{proof}
Supposons que deux points distincts $M$ et $N$ soient tous deux milieux de $[AB]$.
Alors ils appartiennent tous les deux au segment $[AB]$ et vérifient
\[
AM=MB \qquad \text{et} \qquad AN=NB.
\]
Or il ne peut exister qu'un seul point situé exactement à mi-distance de $A$ et de $B$ sur le segment $[AB]$.
Donc $M=N$.
\end{proof}

\begin{exemple}
Si $AB=10$ cm et si $M$ est le milieu de $[AB]$, alors
\[
AM=MB=5 \text{ cm}.
\]
\end{exemple}

%------------------------------------------------
\subsection{Angles, parallélisme et perpendicularité}

\begin{definition}
Un \textbf{angle} est formé par deux demi-droites de même origine.
\end{definition}

\begin{definition}
Un \textbf{angle droit} est un angle mesurant $90^\circ$.
\end{definition}

\begin{definition}
Deux droites sont :
\begin{itemize}
    \item \textbf{perpendiculaires} si elles se coupent en formant un angle droit ;
    \item \textbf{parallèles} si elles ont la même direction.
\end{itemize}
\end{definition}

\begin{propriete}
Si deux droites sont perpendiculaires à une même droite, alors elles sont parallèles.
\end{propriete}

\begin{proof}
Soient $d_1$ et $d_2$ deux droites perpendiculaires à une même droite $d$.
Chacune forme un angle de $90^\circ$ avec $d$.
Elles ont donc la même direction, donc elles sont parallèles.
\end{proof}

\begin{propriete}
Si deux droites sont parallèles, toute droite perpendiculaire à l'une est perpendiculaire à l'autre.
\end{propriete}

\begin{proof}
Deux droites parallèles ont la même direction.
Une droite qui forme un angle droit avec l'une forme donc aussi un angle droit avec l'autre.
\end{proof}

%------------------------------------------------
\subsection{Triangles}

\begin{definition}
Un \textbf{triangle} est une figure formée par trois points non alignés $A$, $B$, $C$ et par les segments $[AB]$, $[BC]$ et $[CA]$.
\end{definition}

\begin{definition}
Un triangle est :
\begin{itemize}
    \item \textbf{isocèle} s'il possède deux côtés de même longueur ;
    \item \textbf{équilatéral} si ses trois côtés sont de même longueur ;
    \item \textbf{rectangle} s'il possède un angle droit.
\end{itemize}
\end{definition}

\begin{theoreme}
Dans tout triangle, la somme des angles vaut $180^\circ$.
\end{theoreme}

\begin{proof}
Soit un triangle $ABC$.
Par le point $A$, traçons la droite $d$ parallèle à $(BC)$.
Les angles formés avec cette droite sont égaux aux angles en $B$ et en $C$ par la propriété des angles alternes-internes.
Or sur une droite, la somme des angles adjacents vaut $180^\circ$.
Donc
\[
\widehat{ABC}+\widehat{BAC}+\widehat{BCA}=180^\circ.
\]
\end{proof}

\begin{corollaire}
Dans un triangle rectangle, la somme des deux angles aigus vaut $90^\circ$.
\end{corollaire}

\begin{propriete}
Dans un triangle isocèle, les angles à la base sont égaux.
\end{propriete}

\begin{proof}
Soit $ABC$ isocèle en $A$, donc $AB=AC$.
Le triangle est symétrique par rapport à la droite passant par $A$ et le milieu de $[BC]$.
Les angles en $B$ et en $C$ sont donc égaux.
\end{proof}

\begin{propriete}
Dans un triangle équilatéral, chaque angle mesure $60^\circ$.
\end{propriete}

\begin{proof}
Les trois angles sont égaux et leur somme vaut $180^\circ$.
Chacun vaut donc
\[
\frac{180^\circ}{3}=60^\circ.
\]
\end{proof}

%------------------------------------------------
\subsection{Médiatrice, bissectrice, hauteur, médiane}

\begin{definition}
La \textbf{médiatrice} d'un segment $[AB]$ est la droite perpendiculaire à $(AB)$ passant par le milieu de $[AB]$.
\end{definition}

\begin{theoreme}
Un point appartient à la médiatrice de $[AB]$ si et seulement s'il est à égale distance de $A$ et de $B$.
\end{theoreme}

\begin{proof}
\textbf{Sens direct :} si $M$ appartient à la médiatrice de $[AB]$, alors la figure est symétrique par rapport à cette médiatrice, donc
\[
MA=MB.
\]

\textbf{Sens réciproque :} si $MA=MB$, alors $M$ est un point équidistant de $A$ et de $B$.
L'ensemble des points équidistants de $A$ et de $B$ est précisément la médiatrice de $[AB]$.
\end{proof}

\begin{corollaire}
Les médiatrices des côtés d'un triangle sont concourantes en un point, centre du cercle circonscrit.
\end{corollaire}

\begin{definition}
La \textbf{bissectrice} d'un angle est la demi-droite qui partage cet angle en deux angles égaux.
\end{definition}

\begin{definition}
Dans un triangle :
\begin{itemize}
    \item une \textbf{hauteur} est une droite passant par un sommet et perpendiculaire au côté opposé ;
    \item une \textbf{médiane} est une droite passant par un sommet et par le milieu du côté opposé.
\end{itemize}
\end{definition}

%------------------------------------------------
\subsection{Cercle}

\begin{definition}
Le \textbf{cercle} de centre $O$ et de rayon $r$ est l'ensemble des points $M$ du plan tels que
\[
OM=r.
\]
\end{definition}

\begin{definition}
Une \textbf{corde} est un segment joignant deux points du cercle.

Un \textbf{diamètre} est une corde passant par le centre.
\end{definition}

\begin{theoreme}[Angle inscrit dans un demi-cercle]
Si $A$, $B$, $C$ sont trois points d'un cercle et si $[AC]$ est un diamètre, alors l'angle $\widehat{ABC}$ est droit.
\end{theoreme}

\begin{proof}
Soit $O$ le centre du cercle.
On a
\[
OA=OB=OC.
\]
Les triangles $AOB$ et $BOC$ sont isocèles.
En additionnant les angles à la base convenablement, on obtient
\[
\widehat{ABC}=90^\circ.
\]
\end{proof}

%------------------------------------------------
\subsection{Transformations du plan}

\begin{definition}
La \textbf{symétrie axiale} d'axe $(d)$ envoie un point $M$ sur un point $M'$ tel que $(d)$ soit la médiatrice de $[MM']$.
\end{definition}

\begin{definition}
La \textbf{symétrie centrale} de centre $O$ envoie un point $M$ sur un point $M'$ tel que $O$ soit le milieu de $[MM']$.
\end{definition}

\begin{definition}
Une \textbf{translation} de vecteur $\uu$ déplace tout point $M$ vers un point $M'$ tel que
\[
\overrightarrow{MM'}=\uu.
\]
\end{definition}

\begin{definition}
Une \textbf{rotation} de centre $O$ et d'angle $\theta$ transforme un point $M$ en un point $M'$ tel que
\[
OM=OM'
\]
et que l'angle orienté de $\overrightarrow{OM}$ vers $\overrightarrow{OM'}$ soit $\theta$.
\end{definition}

\begin{definition}
Une \textbf{homothétie} de centre $O$ et de rapport $k$ transforme un point $M$ en un point $M'$ tel que
\[
\overrightarrow{OM'}=k\,\overrightarrow{OM}.
\]
\end{definition}

\begin{propriete}
Les symétries, translations et rotations conservent les longueurs et les angles.
\end{propriete}

\begin{propriete}
Une homothétie conserve l'alignement et le parallélisme, et multiplie les longueurs par $|k|$.
\end{propriete}

%------------------------------------------------
\subsection{Théorème de Thalès}

\begin{theoreme}[Thalès]
Soient $A$, $B$, $M$ alignés dans cet ordre, et $A$, $C$, $N$ alignés dans cet ordre.
Si
\[
(MN)\parallel(BC),
\]
alors
\[
\frac{AM}{AB}=\frac{AN}{AC}=\frac{MN}{BC}.
\]
\end{theoreme}

\begin{proof}
Les triangles $AMN$ et $ABC$ ont leurs angles correspondants égaux grâce au parallélisme.
Ils sont donc semblables.
Les côtés homologues sont alors proportionnels.
\end{proof}

\begin{theoreme}[Réciproque de Thalès]
Soient $A$, $B$, $M$ alignés et $A$, $C$, $N$ alignés.
Si
\[
\frac{AM}{AB}=\frac{AN}{AC},
\]
alors
\[
(MN)\parallel(BC).
\]
\end{theoreme}

\begin{proof}
On trace par $M$ la parallèle à $(BC)$, qui coupe $(AC)$ en un point $N'$.
Par le théorème de Thalès,
\[
\frac{AM}{AB}=\frac{AN'}{AC}.
\]
Or, par hypothèse,
\[
\frac{AM}{AB}=\frac{AN}{AC}.
\]
Donc $AN'=AN$, d'où $N'=N$.
Ainsi $(MN)\parallel(BC)$.
\end{proof}

%------------------------------------------------
\subsection{Théorème de Pythagore}

\begin{theoreme}[Pythagore]
Si le triangle $ABC$ est rectangle en $A$, alors
\[
BC^2=AB^2+AC^2.
\]
\end{theoreme}

\begin{proof}
Plaçons le triangle dans un repère :
\[
A(0,0),\quad B(b,0),\quad C(0,c).
\]
Alors
\[
AB=b,\qquad AC=c.
\]
La distance $BC$ vaut
\[
BC=\sqrt{(b-0)^2+(0-c)^2}=\sqrt{b^2+c^2}.
\]
Donc
\[
BC^2=b^2+c^2=AB^2+AC^2.
\]
\end{proof}

\begin{theoreme}[Réciproque de Pythagore]
Si dans un triangle $ABC$,
\[
BC^2=AB^2+AC^2,
\]
alors le triangle est rectangle en $A$.
\end{theoreme}

\begin{proof}
On construit un triangle rectangle en $A'$ tel que
\[
A'B=AB,\qquad A'C=AC.
\]
Dans ce triangle,
\[
BC^2=A'B^2+A'C^2=AB^2+AC^2.
\]
Le triangle $ABC$ a donc les mêmes trois côtés que le triangle rectangle construit.
Ils sont isométriques, donc $ABC$ est rectangle en $A$.
\end{proof}

%------------------------------------------------
\subsection{Trigonométrie dans le triangle rectangle}

\begin{definition}
Dans un triangle rectangle, pour un angle aigu $\theta$ :
\[
\cos\theta=\frac{\text{adjacent}}{\text{hypoténuse}},\qquad
\sin\theta=\frac{\text{opposé}}{\text{hypoténuse}},\qquad
\tan\theta=\frac{\text{opposé}}{\text{adjacent}}.
\]
\end{definition}

\begin{propriete}
Pour tout angle aigu $\theta$,
\[
\tan\theta=\frac{\sin\theta}{\cos\theta}.
\]
\end{propriete}

\begin{proof}
On a
\[
\frac{\sin\theta}{\cos\theta}
=
\frac{\frac{\text{opposé}}{\text{hypoténuse}}}{\frac{\text{adjacent}}{\text{hypoténuse}}}
=
\frac{\text{opposé}}{\text{adjacent}}
=
\tan\theta.
\]
\end{proof}

\newpage

%================================================
\section{Seconde : géométrie repérée et vecteurs du plan}
%================================================

\subsection{Repère du plan}

\begin{definition}
Un \textbf{repère} du plan est constitué d'une origine $O$ et de deux axes.

Un point $M$ est repéré par ses coordonnées $(x_M,y_M)$.
\end{definition}

\begin{definition}
Un repère est \textbf{orthonormé} si les axes sont perpendiculaires et si l'unité est la même sur les deux axes.
\end{definition}

%------------------------------------------------
\subsection{Distance et milieu}

\begin{theoreme}
Dans un repère orthonormé, si
\[
A(x_A,y_A),\qquad B(x_B,y_B),
\]
alors
\[
AB=\sqrt{(x_B-x_A)^2+(y_B-y_A)^2}.
\]
\end{theoreme}

\begin{proof}
Le déplacement de $A$ à $B$ se décompose en un déplacement horizontal de longueur $|x_B-x_A|$ et un déplacement vertical de longueur $|y_B-y_A|$.
Ces deux déplacements sont perpendiculaires.
On applique donc le théorème de Pythagore.
\end{proof}

\begin{proposition}
Le milieu du segment $[AB]$ a pour coordonnées
\[
\left(\frac{x_A+x_B}{2},\frac{y_A+y_B}{2}\right).
\]
\end{proposition}

\begin{proof}
Le milieu partage le segment en deux parties égales : ses coordonnées sont donc les moyennes des coordonnées des extrémités.
\end{proof}

%------------------------------------------------
\subsection{Vecteurs du plan}

\begin{definition}
Un \textbf{vecteur} représente un déplacement.
Le vecteur allant de $A$ à $B$ se note $\overrightarrow{AB}$.
\end{definition}

\begin{definition}
Deux vecteurs sont égaux s'ils ont même direction, même sens et même longueur.
\end{definition}

\begin{definition}
Le vecteur nul se note $\overrightarrow{0}$.
\end{definition}

\begin{propriete}[Relation de Chasles]
Pour tous points $A$, $B$, $C$,
\[
\overrightarrow{AB}+\overrightarrow{BC}=\overrightarrow{AC}.
\]
\end{propriete}

\begin{proof}
Aller de $A$ à $B$, puis de $B$ à $C$, revient à aller directement de $A$ à $C$.
\end{proof}

\begin{propriete}
Pour tout point $A$,
\[
\overrightarrow{AA}=\overrightarrow{0},
\qquad
\overrightarrow{AB}=-\overrightarrow{BA}.
\]
\end{propriete}

\begin{proposition}
Si
\[
A(x_A,y_A),\qquad B(x_B,y_B),
\]
alors
\[
\overrightarrow{AB}=(x_B-x_A,\ y_B-y_A).
\]
\end{proposition}

\begin{definition}
La \textbf{norme} d'un vecteur $\uu$ est sa longueur, notée $\norme{\uu}$.
\end{definition}

\begin{propriete}
Si $\uu=(x,y)$, alors
\[
\norme{\uu}=\sqrt{x^2+y^2}.
\]
\end{propriete}

%------------------------------------------------
\subsection{Colinéarité et alignement}

\begin{definition}
Deux vecteurs $\uu$ et $\vv$ sont \textbf{colinéaires} s'il existe un réel $\lambda$ tel que
\[
\uu=\lambda\vv.
\]
\end{definition}

\begin{proposition}
Si
\[
\uu=(x,y),\qquad \vv=(x',y'),
\]
alors $\uu$ et $\vv$ sont colinéaires si et seulement si
\[
xy'-yx'=0.
\]
\end{proposition}

\begin{proof}
Si $\uu=\lambda\vv$, alors
\[
x=\lambda x',\qquad y=\lambda y',
\]
donc
\[
xy'-yx'=\lambda x'y'-\lambda y'x'=0.
\]

Réciproquement, si $xy'-yx'=0$, on montre qu'il existe un réel $\lambda$ tel que $\uu=\lambda\vv$.
\end{proof}

\begin{corollaire}
Les points $A$, $B$, $C$ sont alignés si et seulement si les vecteurs $\overrightarrow{AB}$ et $\overrightarrow{AC}$ sont colinéaires.
\end{corollaire}

%------------------------------------------------
\subsection{Équation de droite}

\begin{definition}
Une droite non verticale du plan admet une équation de la forme
\[
y=mx+p,
\]
où $m$ est le coefficient directeur et $p$ l'ordonnée à l'origine.
\end{definition}

\begin{propriete}
Si $A(x_A,y_A)$ et $B(x_B,y_B)$ avec $x_A\neq x_B$, alors le coefficient directeur de la droite $(AB)$ est
\[
m=\frac{y_B-y_A}{x_B-x_A}.
\]
\end{propriete}

\begin{propriete}
Deux droites non verticales sont parallèles si et seulement si elles ont le même coefficient directeur.
\end{propriete}

%------------------------------------------------
\subsection{Équation d'un cercle}

\begin{definition}
Le cercle de centre $O(x_0,y_0)$ et de rayon $r$ a pour équation
\[
(x-x_0)^2+(y-y_0)^2=r^2.
\]
\end{definition}

\begin{proof}
Un point $M(x,y)$ appartient au cercle si et seulement si
\[
OM=r.
\]
En élevant au carré,
\[
(x-x_0)^2+(y-y_0)^2=r^2.
\]
\end{proof}

\newpage

%================================================
\section{Première spécialité : produit scalaire et géométrie du plan}
%================================================

\subsection{Définition du produit scalaire}

\begin{definition}
Le \textbf{produit scalaire} de deux vecteurs $\uu$ et $\vv$ est défini par
\[
\ps{\uu}{\vv}=\norme{\uu}\,\norme{\vv}\cos\theta,
\]
où $\theta$ est l'angle entre les deux vecteurs.
\end{definition}

\begin{propriete}
Le produit scalaire est symétrique :
\[
\ps{\uu}{\vv}=\ps{\vv}{\uu}.
\]
\end{propriete}

\begin{propriete}
Pour tout vecteur $\uu$,
\[
\ps{\uu}{\uu}=\norme{\uu}^2.
\]
\end{propriete}

\begin{proof}
L'angle entre $\uu$ et lui-même vaut $0$, donc
\[
\ps{\uu}{\uu}=\norme{\uu}\,\norme{\uu}\cos 0=\norme{\uu}^2.
\]
\end{proof}

%------------------------------------------------
\subsection{Orthogonalité}

\begin{theoreme}
Deux vecteurs non nuls sont orthogonaux si et seulement si leur produit scalaire est nul.
\end{theoreme}

\begin{proof}
On a
\[
\ps{\uu}{\vv}=\norme{\uu}\,\norme{\vv}\cos\theta.
\]
Comme les normes sont non nulles, ce produit est nul si et seulement si
\[
\cos\theta=0,
\]
c'est-à-dire si et seulement si $\theta=90^\circ$.
\end{proof}

%------------------------------------------------
\subsection{Expression coordonnée du produit scalaire}

\begin{theoreme}
Dans un repère orthonormé, si
\[
\uu=(x,y),\qquad \vv=(x',y'),
\]
alors
\[
\ps{\uu}{\vv}=xx'+yy'.
\]
\end{theoreme}

\begin{proof}
Dans une base orthonormée $(\overrightarrow{i},\overrightarrow{j})$,
\[
\uu=x\overrightarrow{i}+y\overrightarrow{j},
\qquad
\vv=x'\overrightarrow{i}+y'\overrightarrow{j}.
\]
En développant le produit scalaire, on obtient
\[
\ps{\uu}{\vv}=xx'+yy',
\]
car
\[
\ps{\overrightarrow{i}}{\overrightarrow{i}}=1,\quad
\ps{\overrightarrow{j}}{\overrightarrow{j}}=1,\quad
\ps{\overrightarrow{i}}{\overrightarrow{j}}=0.
\]
\end{proof}

\begin{corollaire}
Dans un repère orthonormé,
\[
\norme{(x,y)}=\sqrt{x^2+y^2}.
\]
\end{corollaire}

\begin{proof}
On écrit
\[
\norme{(x,y)}^2=\ps{(x,y)}{(x,y)}=x^2+y^2.
\]
\end{proof}

%------------------------------------------------
\subsection{Applications}

\begin{propriete}
Deux vecteurs $\uu=(x,y)$ et $\vv=(x',y')$ sont orthogonaux si et seulement si
\[
xx'+yy'=0.
\]
\end{propriete}

\begin{propriete}
Dans un repère orthonormé, deux droites non verticales de coefficients directeurs $m$ et $m'$ sont perpendiculaires si et seulement si
\[
mm'=-1.
\]
\end{propriete}

\begin{proof}
Un vecteur directeur de la première droite est $(1,m)$ et de la seconde $(1,m')$.
La perpendicularité équivaut à
\[
\ps{(1,m)}{(1,m')}=0,
\]
soit
\[
1+mm'=0.
\]
\end{proof}

\begin{theoreme}[Formule d'Al-Kashi]
Dans tout triangle $ABC$,
\[
BC^2=AB^2+AC^2-2AB\cdot AC\cos(\widehat{BAC}).
\]
\end{theoreme}

\begin{proof}
On a
\[
\overrightarrow{BC}=\overrightarrow{BA}+\overrightarrow{AC}.
\]
Donc
\[
BC^2=\ps{\overrightarrow{BC}}{\overrightarrow{BC}}
=\ps{\overrightarrow{BA}+\overrightarrow{AC}}{\overrightarrow{BA}+\overrightarrow{AC}}.
\]
En développant :
\[
BC^2=BA^2+AC^2+2\ps{\overrightarrow{BA}}{\overrightarrow{AC}}.
\]
Or l'angle entre $\overrightarrow{BA}$ et $\overrightarrow{AC}$ vaut $\pi-\widehat{BAC}$, donc
\[
\ps{\overrightarrow{BA}}{\overrightarrow{AC}}
=
-AB\cdot AC\cos(\widehat{BAC}).
\]
D'où
\[
BC^2=AB^2+AC^2-2AB\cdot AC\cos(\widehat{BAC}).
\]
\end{proof}

%------------------------------------------------
\subsection{Centre de gravité}

\begin{theoreme}
Dans un triangle, les trois médianes sont concourantes.
\end{theoreme}

\begin{proof}
Soit $I$ le milieu de $[BC]$.
Alors
\[
\overrightarrow{AI}=\frac12(\overrightarrow{AB}+\overrightarrow{AC}).
\]
Définissons $G$ par
\[
\overrightarrow{AG}=\frac13(\overrightarrow{AB}+\overrightarrow{AC}).
\]
Alors
\[
\overrightarrow{AG}=\frac23\,\overrightarrow{AI},
\]
donc $G$ appartient à la médiane issue de $A$.
De même, on montre qu'il appartient aux deux autres médianes.
\end{proof}

\begin{corollaire}
Le centre de gravité partage chaque médiane dans le rapport
\[
AG=\frac23\,AI.
\]
\end{corollaire}

%------------------------------------------------
\subsection{Cercle trigonométrique}

\begin{definition}
Le \textbf{cercle trigonométrique} est le cercle de centre $O$ et de rayon $1$ dans un repère orthonormé.
\end{definition}

\begin{definition}
À tout réel $x$, on associe sur le cercle trigonométrique le point de coordonnées
\[
(\cos x,\sin x).
\]
\end{definition}

\begin{propriete}
Pour tout réel $x$,
\[
\cos^2 x+\sin^2 x=1.
\]
\end{propriete}

\begin{proof}
Le point de coordonnées $(\cos x,\sin x)$ appartient au cercle trigonométrique, donc sa distance à l'origine vaut $1$.
\end{proof}

\newpage

%================================================
\section{Terminale spécialité : géométrie dans l'espace}
%================================================

\subsection{Repère de l'espace}

\begin{definition}
Un repère de l'espace est constitué d'une origine $O$ et de trois axes non coplanaires.
Un point $M$ est repéré par ses coordonnées
\[
M(x,y,z).
\]
\end{definition}

\begin{definition}
Un repère est \textbf{orthonormé} si les trois axes sont deux à deux perpendiculaires et si l'unité est la même sur chacun.
\end{definition}

%------------------------------------------------
\subsection{Vecteurs et distance dans l'espace}

\begin{proposition}
Si
\[
A(x_A,y_A,z_A),\qquad B(x_B,y_B,z_B),
\]
alors
\[
\overrightarrow{AB}=(x_B-x_A,\ y_B-y_A,\ z_B-z_A).
\]
\end{proposition}

\begin{propriete}
Pour tous points $A$, $B$, $C$,
\[
\overrightarrow{AB}+\overrightarrow{BC}=\overrightarrow{AC}.
\]
\end{propriete}

\begin{theoreme}
Dans un repère orthonormé de l'espace,
\[
AB=\sqrt{(x_B-x_A)^2+(y_B-y_A)^2+(z_B-z_A)^2}.
\]
\end{theoreme}

\begin{proof}
On applique la formule de la norme au vecteur $\overrightarrow{AB}$.
\end{proof}

%------------------------------------------------
\subsection{Produit scalaire dans l'espace}

\begin{theoreme}
Si
\[
\uu=(x,y,z),\qquad \vv=(x',y',z'),
\]
dans une base orthonormée de l'espace, alors
\[
\ps{\uu}{\vv}=xx'+yy'+zz'.
\]
\end{theoreme}

\begin{corollaire}
On a
\[
\norme{(x,y,z)}=\sqrt{x^2+y^2+z^2}.
\]
\end{corollaire}

\begin{propriete}
Deux vecteurs non nuls de l'espace sont orthogonaux si et seulement si leur produit scalaire est nul.
\end{propriete}

%------------------------------------------------
\subsection{Droites de l'espace}

\begin{definition}
Une droite passant par $A(x_A,y_A,z_A)$ et de vecteur directeur $\uu=(a,b,c)$ admet une représentation paramétrique
\[
\left\{
\begin{array}{l}
x=x_A+ta\\
y=y_A+tb\\
z=z_A+tc
\end{array}
\right.
\qquad (t\in\R).
\]
\end{definition}

\begin{propriete}
Deux droites sont parallèles si et seulement si leurs vecteurs directeurs sont colinéaires.
\end{propriete}

%------------------------------------------------
\subsection{Plans de l'espace}

\begin{definition}
Un plan de l'espace admet une équation cartésienne de la forme
\[
ax+by+cz+d=0,
\]
avec
\[
(a,b,c)\neq(0,0,0).
\]
\end{definition}

\begin{definition}
Le vecteur
\[
\overrightarrow{n}=(a,b,c)
\]
est appelé \textbf{vecteur normal} au plan.
\end{definition}

\begin{proposition}
Le plan passant par $A(x_A,y_A,z_A)$ et de vecteur normal $\overrightarrow{n}=(a,b,c)$ a pour équation
\[
a(x-x_A)+b(y-y_A)+c(z-z_A)=0.
\]
\end{proposition}

\begin{proof}
Un point $M(x,y,z)$ appartient au plan si et seulement si le vecteur $\overrightarrow{AM}$ est orthogonal au vecteur normal :
\[
\ps{\overrightarrow{AM}}{\overrightarrow{n}}=0.
\]
Cela donne
\[
a(x-x_A)+b(y-y_A)+c(z-z_A)=0.
\]
\end{proof}

\begin{propriete}
Deux plans sont parallèles si et seulement si leurs vecteurs normaux sont colinéaires.
\end{propriete}

%------------------------------------------------
\subsection{Orthogonalité dans l'espace}

\begin{definition}
Une droite est perpendiculaire à un plan si elle est perpendiculaire à toute droite du plan passant par son pied.
\end{definition}

\begin{propriete}
Une droite de vecteur directeur $\uu$ est perpendiculaire à un plan de vecteur normal $\overrightarrow{n}$ si et seulement si $\uu$ et $\overrightarrow{n}$ sont colinéaires.
\end{propriete}

\begin{propriete}
Une droite de vecteur directeur $\uu$ est parallèle à un plan de vecteur normal $\overrightarrow{n}$ si et seulement si
\[
\ps{\uu}{\overrightarrow{n}}=0.
\]
\end{propriete}

%------------------------------------------------
\subsection{Sphère}

\begin{definition}
Une \textbf{sphère} de centre $O(x_0,y_0,z_0)$ et de rayon $r$ est l'ensemble des points $M(x,y,z)$ tels que
\[
OM=r.
\]
\end{definition}

\begin{propriete}
L'équation cartésienne d'une sphère de centre $O(x_0,y_0,z_0)$ et de rayon $r$ est
\[
(x-x_0)^2+(y-y_0)^2+(z-z_0)^2=r^2.
\]
\end{propriete}

\begin{proof}
C'est la traduction analytique de la condition
\[
OM=r
\]
à l'aide de la formule de distance dans l'espace.
\end{proof}

%------------------------------------------------
\subsection{Projeté orthogonal}

\begin{definition}
Le \textbf{projeté orthogonal} d'un point $M$ sur un plan $\mathcal P$ est le point $H$ du plan tel que la droite $(MH)$ soit perpendiculaire au plan.
\end{definition}

\begin{theoreme}
Le projeté orthogonal d'un point sur un plan est le point du plan le plus proche de ce point.
\end{theoreme}

\begin{proof}
Soit $H$ le projeté orthogonal de $M$ sur le plan $\mathcal P$, et soit $N$ un point quelconque du plan.
Le triangle $MHN$ est rectangle en $H$.
Par le théorème de Pythagore,
\[
MN^2=MH^2+HN^2\ge MH^2.
\]
Donc
\[
MN\ge MH.
\]
Ainsi, $H$ est le point du plan le plus proche de $M$.
\end{proof}

\newpage

%================================================
\section{Méthodes classiques de démonstration}
%================================================

\begin{methode}
Pour montrer que des points sont alignés :
\begin{itemize}
    \item exhiber une droite commune ;
    \item montrer que deux vecteurs sont colinéaires ;
    \item montrer qu'ils vérifient une même équation de droite.
\end{itemize}
\end{methode}

\begin{methode}
Pour montrer que deux droites sont parallèles :
\begin{itemize}
    \item utiliser les propriétés d'angles avec des parallèles ;
    \item utiliser la réciproque de Thalès ;
    \item montrer que des vecteurs directeurs sont colinéaires ;
    \item comparer les coefficients directeurs.
\end{itemize}
\end{methode}

\begin{methode}
Pour montrer qu'il y a perpendicularité :
\begin{itemize}
    \item repérer un angle droit ;
    \item utiliser Pythagore ou sa réciproque ;
    \item montrer que le produit scalaire est nul.
\end{itemize}
\end{methode}

\begin{methode}
Pour calculer une longueur :
\begin{itemize}
    \item utiliser les égalités géométriques connues ;
    \item appliquer Pythagore ;
    \item utiliser la formule de distance ;
    \item utiliser le produit scalaire ou Al-Kashi.
\end{itemize}
\end{methode}

%================================================
\section{Exemples de synthèse}
%================================================

\begin{exemple}[Alignement dans le plan]
Soient
\[
A(1,2),\quad B(3,5),\quad C(5,8).
\]
On calcule
\[
\overrightarrow{AB}=(2,3),\qquad \overrightarrow{AC}=(4,6)=2(2,3).
\]
Les vecteurs $\overrightarrow{AB}$ et $\overrightarrow{AC}$ sont colinéaires.
Donc les points $A$, $B$, $C$ sont alignés.
\end{exemple}

\begin{exemple}[Milieu]
Soient
\[
A(-1,4),\qquad B(5,2).
\]
Le milieu de $[AB]$ est
\[
M\left(\frac{-1+5}{2},\frac{4+2}{2}\right)=(2,3).
\]
\end{exemple}

\begin{exemple}[Orthogonalité dans le plan]
Soient
\[
\uu=(2,-3),\qquad \vv=(3,2).
\]
Alors
\[
\ps{\uu}{\vv}=2\cdot 3+(-3)\cdot 2=6-6=0.
\]
Donc les vecteurs sont orthogonaux.
\end{exemple}

\begin{exemple}[Droite dans l'espace]
La droite passant par
\[
A(1,0,2)
\]
et de vecteur directeur
\[
(2,-1,3)
\]
admet pour représentation paramétrique
\[
\left\{
\begin{array}{l}
x=1+2t\\
y=-t\\
z=2+3t
\end{array}
\right.
\]
\end{exemple}

\begin{exemple}[Plan dans l'espace]
Déterminons une équation du plan passant par
\[
A(1,2,3)
\]
et de vecteur normal
\[
(2,-1,4).
\]
On écrit
\[
2(x-1)-(y-2)+4(z-3)=0.
\]
Après développement :
\[
2x-y+4z-8=0.
\]
\end{exemple}

%================================================
\section{Formulaire récapitulatif}
%================================================

\subsection*{Dans le plan}
\[
AB=\sqrt{(x_B-x_A)^2+(y_B-y_A)^2}
\]
\[
M_{[AB]}=\left(\frac{x_A+x_B}{2},\frac{y_A+y_B}{2}\right)
\]
\[
\overrightarrow{AB}=(x_B-x_A,\ y_B-y_A)
\]
\[
\ps{(x,y)}{(x',y')}=xx'+yy'
\]
\[
\norme{(x,y)}=\sqrt{x^2+y^2}
\]
\[
\text{colinéarité : } xy'-yx'=0
\]
\[
\text{droite non verticale : } y=mx+p
\]
\[
\text{cercle : } (x-x_0)^2+(y-y_0)^2=r^2
\]

\subsection*{Dans l'espace}
\[
AB=\sqrt{(x_B-x_A)^2+(y_B-y_A)^2+(z_B-z_A)^2}
\]
\[
\overrightarrow{AB}=(x_B-x_A,\ y_B-y_A,\ z_B-z_A)
\]
\[
\ps{(x,y,z)}{(x',y',z')}=xx'+yy'+zz'
\]
\[
\norme{(x,y,z)}=\sqrt{x^2+y^2+z^2}
\]
\[
\text{droite : }
\left\{
\begin{array}{l}
x=x_A+ta\\
y=y_A+tb\\
z=z_A+tc
\end{array}
\right.
\]
\[
\text{plan : } ax+by+cz+d=0
\]
\[
\text{sphère : } (x-x_0)^2+(y-y_0)^2+(z-z_0)^2=r^2
\]

%================================================
\section*{Conclusion}
%================================================

Ce cours regroupe l'essentiel de la géométrie jusqu'à la terminale :
\begin{itemize}
    \item géométrie plane classique ;
    \item transformations ;
    \item Thalès, Pythagore, trigonométrie ;
    \item géométrie repérée et vectorielle ;
    \item produit scalaire ;
    \item géométrie de l'espace.
\end{itemize}

Pour bien réviser, il faut :
\begin{itemize}
    \item apprendre précisément les définitions ;
    \item connaître les hypothèses des théorèmes ;
    \item refaire les démonstrations ;
    \item s'entraîner sur des exercices variés.
\end{itemize}

\end{document}